\section{Testing strategy}

\subsection{Required test categories}

The existing suite checks driver behavior, PV naming and aggregation, update-period limits, XML validity of generated \file{.bob} resources, complete PV coverage, and controls for writable PVs. Extend the suite when changing any of those contracts.

For a new feature, include as applicable:

\begin{itemize}
  \item pure parsing and command-format tests;
  \item driver tests with fake transports;
  \item IOC putter and polling tests;
  \item configuration validation tests;
  \item duplicate-PV and profile-composition tests;
  \item generated-display contract tests;
  \item operator script tests with fake processes or commands;
  \item a documented manual hardware acceptance test.
\end{itemize}

\subsection{Pre-commit checklist}

\begin{lstlisting}[style=shell]
source .venv/bin/activate
ruff check src tests
pytest
bdx-pv-list --config-dir config/profiles/default > /tmp/bdx-default.pvs
bdx-generate-displays --config-dir config/profiles/default \
  --output-dir phoebus/displays

git status --short
git diff --check
\end{lstlisting}

Inspect generated XML diffs rather than committing them blindly. Confirm that no local runtime file, credentials, logs, databases, or archive data are staged.

\section{Logging and diagnostics}

The package configures Python logging centrally. Set \code{BDX_LOG_LEVEL} to increase verbosity during diagnosis:

\begin{lstlisting}[style=shell]
BDX_LOG_LEVEL=DEBUG bdx-prototype-ioc \
  --config-dir config/profiles/default
\end{lstlisting}

Useful diagnostic principles:

\begin{itemize}
  \item start with read-only network checks;
  \item distinguish a failed TCP connection from a protocol parse failure;
  \item report the device, channel, attempted operation, and underlying exception;
  \item avoid repeated identical log floods during a persistent poll failure;
  \item log recovery when communication becomes healthy again;
  \item never print credentials or private keys.
\end{itemize}

For Channel Access checks, set an explicit address list and disable automatic broadcast discovery when diagnosing a split-host system:

\begin{lstlisting}[style=shell]
export EPICS_CA_ADDR_LIST="172.22.50.2 172.22.50.10"
export EPICS_CA_AUTO_ADDR_LIST=NO
caproto-get BDX:GLOBAL:SYSTEM_STATE
caproto-get BDX:ENV:TEMP:T00:VALUE
\end{lstlisting}

\section{Git workflow and stable releases}

Routine development should occur on \code{devel}. Use focused commits with tests and documentation in the same change. Promote to \code{main} only after the intended hardware/simulation workflow and generated resources have been validated.

Recommended workflow:

\begin{lstlisting}[style=shell]
git switch devel
git pull --ff-only

# edit, test, regenerate, document

git status --short
git diff --check
git add <reviewed paths>
git commit -m "Describe the focused change"
git push origin devel
\end{lstlisting}

For a stable promotion:

\begin{enumerate}
  \item ensure \code{devel} is clean and pushed;
  \item run the complete automated test suite;
  \item perform the relevant hardware acceptance checks;
  \item verify generated Phoebus resources and Archiver lists;
  \item update manuals and version/revision notes;
  \item build and inspect the two PDF manuals;
  \item merge or fast-forward the validated state to \code{main} according to the chosen repository policy;
  \item tag the release if a durable deployment reference is required.
\end{enumerate}

Avoid committing directly to \code{main} for ordinary development.

\section{Building these manuals}

The LaTeX sources are in \file{documentation}. From the repository root:

\begin{lstlisting}[style=shell]
make -C documentation
\end{lstlisting}

The build places the final operator and developer PDFs in the repository root:

\begin{lstlisting}
BDX_Slow_Control_Quickstart.pdf
BDX_Slow_Control_Developer_Manual.pdf
\end{lstlisting}

Intermediate LaTeX files remain under \file{documentation/build} and are ignored or removed with:

\begin{lstlisting}[style=shell]
make -C documentation clean
\end{lstlisting}

Before committing exported PDFs, render and inspect every page. Check for clipped code, overfull lines, broken links, stale addresses, and inconsistency with the current branch.

\section{Change checklist}

Use this compact checklist for modifications that affect operation or the PV contract.

\begin{enumerate}
  \item Update or add the driver interface, simulation, hardware backend, IOC, builder, factory, CLI, and profile as required.
  \item Preserve request/setpoint/readback semantics and common health PV behavior.
  \item Add or update automated tests.
  \item Regenerate only the affected Phoebus resources, then inspect all diffs.
  \item Update Archiver PV lists for newly archived PVs.
  \item Update the quickstart guide if operator behavior changes.
  \item Update this manual and focused deployment documents if architecture or maintenance changes.
  \item Run Ruff, pytest, PV-list generation, and display generation.
  \item Perform a controlled hardware acceptance test when hardware behavior changed.
  \item Commit to \code{devel}; promote to \code{main} only after validation.
\end{enumerate}
